\section{Technical Appendix}
\label{sec:appendix}

\subsection{A. The architectural rate bound}
\paragraph{Proposition 1 (global ceiling).}
For the head $g=o\odot\tanh(h_t)$, $z_{\mathrm{raw}}=w^{\!\top}g+b$,
$z=z_{\max}-\mathrm{softplus}(z_{\max}-z_{\mathrm{raw}})$,
$\lambda=\varsigma\,\mathrm{softplus}(z)$ with $\varsigma>0$:
\[
\lambda(t)\;\le\;\bar\lambda:=\varsigma\,\mathrm{softplus}(z_{\max})
\]
for every $h_t$ and every $(W_o,b_o,w,b)$.
\emph{Proof.} $\mathrm{softplus}(x)>0$ for all $x$, so
$z=z_{\max}-\mathrm{softplus}(z_{\max}-z_{\mathrm{raw}})<z_{\max}$;
$\mathrm{softplus}$ is increasing, so
$\lambda<\varsigma\,\mathrm{softplus}(z_{\max})$. \hfill$\square$

\paragraph{Consequences.} (i) The conditional intensity is globally bounded,
so for any finite horizon $T$ the compensator satisfies
$\Lambda(T)\le\bar\lambda T<\infty$ and the process is non-explosive on
$[0,T]$. (ii) The event count $N(T)$ is stochastically dominated by a
homogeneous Poisson count with mean $\bar\lambda T$: the process can be
constructed by thinning a homogeneous Poisson process of rate $\bar\lambda$,
and thinning only removes points. (iii) $\bar\lambda$ is an exact dominating
rate for Ogata thinning \citep{lewis1979thinning,ogata1981lewis}. The bound
establishes nothing beyond (i)--(iii); stationarity, ergodicity, and
realism are empirical matters.

\paragraph{Proposition 2 (parameter-dependent floor and pre-cap bound).}
Since $g\in(-1,1)^H$, for fixed trained parameters
$b-\lVert w\rVert_1<z_{\mathrm{raw}}<b+\lVert w\rVert_1$. Hence, even
without the cap, $\lambda\le\varsigma\,\mathrm{softplus}(b+\lVert
w\rVert_1)$ is a computable post-training ceiling; and with the cap,
\[
\lambda\;\ge\;\varsigma\,\mathrm{softplus}\bigl(z_{\max}
-\mathrm{softplus}(z_{\max}-b+\lVert w\rVert_1)\bigr)\;>\;0,
\]
a strictly positive, typically numerically negligible floor. The intensity
therefore has zero infimum only in the limit of unbounded parameters, not
for any fixed trained model. The cap's contribution over the pre-cap bound
is that $\bar\lambda$ is user-chosen, independent of the trained magnitudes
of $w$ and $b$, and invariant under training.

\subsection{B. Calibration protocol}
Targets are validation-split rates ($38.27$, $47.92$, $23.98$ ev/s for BTC,
ETH, SOL). For each checkpoint: (1)~probe roll-outs of $600$\,s at
multiplier $\kappa$ starting from a geometric grid until the target is
bracketed; (2)~bisection in $\log\kappa$, at most ten steps, until a probe
lands within $5\%$ of the target; (3)~a full-scale validation roll-out must
reproduce the target within $15\%$; (4)~failures escalate to full-scale
probes, and unresolved failures exclude the checkpoint and are reported.
Probe seeds are deterministic; retried failures reproduce their bisection
brackets bit-identically. Monotonicity of the closed-loop rate in $\kappa$
is not assumed or proven; it held empirically over every successful bracket.
Calibrated multipliers $\kappa$ (seeds 1/2/3): \ssppp{} BTC
$0.545/0.569/0.648$, ETH $0.723/0.707/0.582$, SOL $0.557/0.595/0.621$;
\spppu{} (passing checkpoints only) BTC $0.534$ (s2), $0.751$ (s3), ETH
$1.102$ (s3).

\subsection{C. Metric definitions}
\textbf{NLL} is per-event negative log-likelihood with the compensator
integrated by a 32-point log-spaced quadrature per inter-event gap
(horizon $60$\,s; beyond the grid the intensity is extended at its terminal
value). \textbf{Type accuracy} is top-1 accuracy of $p(e\mid t^{*})$
evaluated at the \emph{observed} next-event time $t^{*}$ (accuracy
conditional on the observed time). \textbf{Expected-time MAE} computes
$\hat{\mathbb{E}}[\Delta t]$ by survival quadrature on the same log-spaced
grid, adds the tail term $S(T_{\max})/\lambda(T_{\max})$, clamps the result
to $[10^{-6},2T_{\max}]$ with $T_{\max}=60$\,s, and reports the mean
absolute error against the true gap. All prediction metrics are computed by
a streaming evaluator with carried state, one cold start per genuine stream
segment, which hard-fails if any test event goes unscored.

\subsection{D. Stylized-fact construction (order-flow proxy)}
The models emit event times and types only, so the battery operates on an
order-flow proxy applied identically to real and simulated streams. Events
are bucketed at $1$\,s. Per bucket, \emph{activity} $a_t$ is the total
event count and the \emph{proxy return} $r_t$ is the signed count of
price-moving events (unit tick): market orders and best-level inside-spread
improvements, bid side positive, ask side negative. F0 is the event rate;
F2 is the Fano factor $\mathrm{Var}/\mathrm{Mean}$ of bucket counts
aggregated at scales $\{1,2,5,10,20,50\}$\,s, averaged over scales; F1, F3--F9,
and F11 are standard stylized-fact statistics of $r_t$ (autocorrelation,
kurtosis, Hill index, skewness, clustering, aggregated kurtosis, power-law
decay of $|r|$ autocorrelation, leverage, timescale asymmetry); F10
correlates $a_t$ with $|r_t|$ (activity--volatility). None of these are
statistics of asset price returns; we label them accordingly.

\subsection{E. Model and training configuration}
All models share the input embedding (channel embedding 64, time embedding
64) and differ only in the decoder. Backbone/decoder hidden size 64, two
layers. Training: AdamW, learning rate $2\times10^{-3}$, weight decay
$10^{-6}$, gradient clipping at norm $0.5$, ReduceLROnPlateau (factor 0.5,
patience 5), batch 64 lanes, window 1024 events with stride 1024
(non-overlapping, TBPTT with detached carried state), 12 epochs, best
checkpoint by validation loss, three training seeds $\{1,2,3\}$, TF32
matmuls, no dropout. \ssppp{} rate head: $z_{\max}=6$; $\varsigma$
log-parameterized and initialized from the training-split rate
($\varsigma_0=\text{rate}/\ln 2$); gate bias and rate bias initialized to
zero; rate weights Xavier-uniform (gain $0.5$). Simulation: carried-state
roll-outs, $32$ sequences $\times$ $600$\,s, three roll-out seeds, thinning
for \ssppp{} (constant ceiling $\varsigma\,\mathrm{softplus}(z_{\max})$),
compensator-inversion sampling for baselines.

\subsection{F. Efficiency table}
\begin{table}[t]
\centering
\small
\setlength{\tabcolsep}{4pt}
\begin{tabular}{lccc}
\toprule
 & train step (s) & vs.\ seq.\ loop & sim.\ ev/s\\
\midrule
\nhp{}    & 4.35  & --           & 1253\\
LSTM      & 0.019 & --           & 2083\\
SAHP      & 0.059 & --           & 583\\
PCT-LSTM  & 0.79  & --           & 1919\\
\spppu{}   & 0.055 & $64\times$   & 1163\\
\ssppp{}  & 0.059 & $60\times$   & 879 / 353$^{\dag}$\\
\bottomrule
\end{tabular}
\caption{Efficiency on one GPU. Training: one fwd+bwd on $64\times1024$
events (scan for the \sppp{} family; speedup vs.\ their sequential loop).
Simulation: events per wall-second, $8\times60$\,s roll-outs.
$^{\dag}$inversion~/ exact-ceiling Ogata thinning.}
\label{tab:speed}
\end{table}

\subsection{G. Dataset}
Coinbase BTC-USD, ETH-USD, and SOL-USD trades and Level-2 book updates,
seven days per asset, encoded into $|\mathcal{E}|=62$ atomic types
($20+20+2+20$ over LO/CO/MO/IS $\times$ side $\times$ level). Genuine test
events scored: $3{,}226{,}821$ (BTC), $3{,}623{,}575$ (ETH), $2{,}000{,}638$
(SOL). Trades are attributed to book changes by inter-update timestamp
interval and netted against liquidity changes at the traded price; residual
decreases are cancellations. Multiple atomic events from one update share
its timestamp and are ordered deterministically (MO, CO, LO, IS; bid before
ask; then by level). Order sizes are retained in the dataset but unused by
the models in this paper.
% AUTHOR ACTION REQUIRED: add exact collection dates and timezone,
% train/val/test boundaries and per-split event counts, number of genuine
% segments per asset, timestamp resolution, and data-quality counters
% (unmatched trades, ambiguous matches, dropped updates).
